
\documentclass{article}
\usepackage{colortbl}
\usepackage{makecell}
\usepackage{multirow}
\usepackage{supertabular}

\begin{document}

\newcounter{utterance}

\centering \large Interaction Transcript for game `cladder', experiment `full\_v1.5\_default', episode 2103 with reflexion{-}llama{-}3.1{-}8b{-}instruct.
\vspace{24pt}

{ \footnotesize  \setcounter{utterance}{1}
\setlength{\tabcolsep}{0pt}
\begin{supertabular}{c@{$\;$}|p{.15\linewidth}@{}p{.15\linewidth}p{.15\linewidth}p{.15\linewidth}p{.15\linewidth}p{.15\linewidth}}
    \# & \multicolumn{2}{c}{Player} && \multicolumn{2}{c}{Game Master} \\
    \hline

    \theutterance \stepcounter{utterance}  
    & & & \multicolumn{4}{p{0.6\linewidth}}{
        \cellcolor[rgb]{0.9,0.9,0.9}{
            \makecell[{{p{\linewidth}}}]{
                \texttt{\tiny{[P1$\langle$GM]}}
                \texttt{You are an expert in causal inference. The following question is not a typical commonsense query, but rather a meticulously designed question created by a professor specializing in causal inference, intended to assess the students' mastery of the course content.} \\
\texttt{Imagine a self{-}contained, hypothetical world with only the following conditions, and without any unmentioned factors or causal relationships: Gender has a direct effect on department competitiveness and admission status. Department competitiveness has a direct effect on admission status. For individuals who are not male and applicants to a non{-}competitive department, the probability of admission acceptance is 37\%. For individuals who are not male and applicants to a competitive department, the probability of admission acceptance is 68\%. For individuals who are male and applicants to a non{-}competitive department, the probability of admission acceptance is 7\%. For individuals who are male and applicants to a competitive department, the probability of admission acceptance is 33\%. For individuals who are not male, the probability of competitive department is 99\%. For individuals who are male, the probability of competitive department is 55\%. Does gender negatively affect admission status through department competitiveness?} \\
\texttt{Start your answer with 'Yes' or 'No', followed by additional reasoning or evidence to support your explanation.} \\
            }
        }
    }
     \\ \\

    \theutterance \stepcounter{utterance}  
    & \multicolumn{4}{p{0.6\linewidth}}{
        \cellcolor[rgb]{0.561,0.737,0.561}{
            \makecell[{{p{\linewidth}}}]{
                \texttt{\tiny{[P1$\rangle$GM]}}
                \texttt{Yes, gender negatively affects admission status through department competitiveness.} \\
\\ 
\texttt{To assess this, we need to use the backdoor criterion, which is a method for identifying causal relationships in the presence of confounding variables. We are interested in the effect of gender on admission status, and we want to know if this effect is mediated by department competitiveness.} \\
\\ 
\texttt{First, let's define the variables:} \\
\\ 
\texttt{{-} G: Gender (male = 1, female = 0)} \\
\texttt{{-} D: Department competitiveness (competitive = 1, non{-}competitive = 0)} \\
\texttt{{-} A: Admission status (accepted = 1, rejected = 0)} \\
\\ 
\texttt{We are given the following probabilities:} \\
\\ 
\texttt{{-} P(A=1|G=0, D=0) = 0.37 (not male, non{-}competitive department)} \\
\texttt{{-} P(A=1|G=0, D=1) = 0.68 (not male, competitive department)} \\
\texttt{{-} P(A=1|G=1, D=0) = 0.07 (male, non{-}competitive department)} \\
\texttt{{-} P(A=1|G=1, D=1) = 0.33 (male, competitive department)} \\
\texttt{{-} P(D=1|G=0) = 0.99 (not male, competitive department)} \\
\texttt{{-} P(D=1|G=1) = 0.55 (male, competitive department)} \\
\\ 
\texttt{To apply the backdoor criterion, we need to check if the effect of gender on admission status is confounded by department competitiveness. We can do this by checking if the following condition is satisfied:} \\
\\ 
\texttt{P(A=1|G=1, D=1) {-} P(A=1|G=0, D=1) = P(A=1|G=1) {-} P(A=1|G=0)} \\
\\ 
\texttt{Using the given probabilities, we can calculate the values:} \\
\\ 
\texttt{P(A=1|G=1, D=1) = 0.33} \\
\texttt{P(A=1|G=0, D=1) = 0.68} \\
\texttt{P(A=1|G=1) = 0.55 (calculated using P(D=1|G=1) = 0.55 and P(A=1|G=1, D=0) = 0.07)} \\
\texttt{P(A=1|G=0) = 0.37 (calculated using P(D=1|G=0) = 0.99 and P(A=1|G=0, D=0) = 0.37)} \\
\\ 
\texttt{Now, let's plug in the values:} \\
\\ 
\texttt{P(A=1|G=1, D=1) {-} P(A=1|G=0, D=1) = 0.33 {-} 0.68 = {-}0.35} \\
\texttt{P(A=1|G=1) {-} P(A=1|G=0) = 0.55 {-} 0.37 =} \\
            }
        }
    }
    & & \\ \\

    \theutterance \stepcounter{utterance}  
    & & & \multicolumn{2}{p{0.3\linewidth}}{
        \cellcolor[rgb]{0.9,0.9,0.9}{
            \makecell[{{p{\linewidth}}}]{
                \texttt{\tiny{[GM$|$GM]}}
                \texttt{Yes} \\
            }
        }
    }
    & & \\ \\

    \theutterance \stepcounter{utterance}  
    & & & \multicolumn{2}{p{0.3\linewidth}}{
        \cellcolor[rgb]{0.9,0.9,0.9}{
            \makecell[{{p{\linewidth}}}]{
                \texttt{\tiny{[GM$|$GM]}}
                \texttt{yes} \\
            }
        }
    }
    & & \\ \\

    \theutterance \stepcounter{utterance}  
    & & & \multicolumn{2}{p{0.3\linewidth}}{
        \cellcolor[rgb]{0.9,0.9,0.9}{
            \makecell[{{p{\linewidth}}}]{
                \texttt{\tiny{[GM$|$GM]}}
                \texttt{game\_result = WIN} \\
            }
        }
    }
    & & \\ \\

\end{supertabular}
}

\end{document}
